\documentclass[10pt]{beamer}
% \usepackage{inputenc}
% \usepackage{graphicx}
% \usepackage{enumerate}
% \usepackage{listings}
\usepackage{caption, subcaption}
\captionsetup[subfigure]{justification=centering}
\captionsetup{font=scriptsize,labelfont=scriptsize}
\usepackage{amsmath}
\usepackage{amssymb}
\usepackage{amsthm}
% \usetheme{metropolis}    
% \usepackage{verbatim}
% \usepackage{alltt}
% \usepackage{xcolor}
% \usepackage{appendixnumberbeamer}
\usepackage{animate}
% \usepackage{xmpmulti}
\usepackage{threeparttable, longtable}
\usepackage{booktabs}
% \usepackage{setspace}
% \usepackage{adjustbox}
% \usepackage[T1]{fontenc}
% \usepackage[colorlinks=true, allcolors=blue]{hyperref}
% % \usepackage{color,soul}
%\usepackage{tikz}
% \usepackage{multirow}
% \usepackage{float}
% \usepackage{placeins}

\usepackage{tikz}
\usetikzlibrary{tikzmark,overlay-beamer-styles, calc}

% \def\signed #1{{\leavevmode\unskip\nobreak\hfil\penalty50\hskip1em
% 		\hbox{}\nobreak\hfill #1%
% 		\parfillskip=0pt \finalhyphendemerits=0 \endgraf}}

% \newsavebox\mybox
% \newenvironment{aquote}[1]
% {\savebox\mybox{#1}\begin{quote}\openautoquote\hspace*{-.7ex}}
% 	{\unskip\closeautoquote\vspace*{1mm}\signed{\usebox\mybox}\end{quote}}
% \newcommand{\backupbegin}{
% 	\newcounter{finalframe}
% 	\setcounter{finalframe}{\value{framenumber}}
% }
% \newcommand{\backupend}{
% 	\setcounter{framenumber}{\value{finalframe}}
% }
% \usepackage[normalem]{ulem}
% \useunder{\uline}{\ul}{}
% \makeatletter
% \setbeamertemplate{footline}
% {
% 	\leavevmode%
% 	\hbox{%
% 		\begin{beamercolorbox}[wd=.15\paperwidth,ht=2.25ex,dp=1ex,center]{institute in head/foot}%
% 			\usebeamerfont{title in head/foot}%
% 			%  \insertshortauthor 
% 		\end{beamercolorbox}%
% 		\begin{beamercolorbox}[wd=.6\paperwidth,ht=2.25ex,dp=1ex,center]{institute in head/foot}%
% 			\usebeamerfont{title in head/foot}%
% 			\insertshorttitle
% 		\end{beamercolorbox}%
% 		\begin{beamercolorbox}[wd=.15\paperwidth,ht=2.25ex,dp=1ex,center]{institute in head/foot}%
% 			\usebeamerfont{title in head/foot}%
% 			\insertshortdate
% 		\end{beamercolorbox}%
% 		\begin{beamercolorbox}[wd=.1\paperwidth,ht=2.25ex,dp=1ex,right]{institute in head/foot}%
% 			\usebeamerfont{title in head/foot} 
% 			% \insertframenumber{} / \inserttotalframenumber\hspace*{2ex} 
%                 \insertframenumber{} \hspace*{2ex} 
% 	\end{beamercolorbox}}%
% }

% \definecolor{blueucl}{RGB}{0,87,156} 
% \setbeamercolor{itemize item}{fg = blueucl}
% \setbeamercolor{itemize subitem}{fg = blueucl}
% \setbeamercolor{enumerate item}{fg = blueucl}
% \setbeamercolor{enumerate subitem}{fg = blueucl}
% \setbeamertemplate{itemize item}[triangle]
% \setbeamercolor{progress bar}{fg=blueucl}
% \setbeamercolor{alerted text}{fg = blueucl}

% \lstset{numbers=left, numberstyle=\tiny, stepnumber=1,firstnumber=1,
% 	numbersep=5pt,
% 	stringstyle=\ttfamily,
% 	basicstyle=\footnotesize, 
% 	showstringspaces=false, 
% 	stringstyle=\color{blueucl}, 
% 	tabsize=1
% }
% \makeatother
% \setbeamertemplate{caption}[numbered]
% \newcommand{\ep}{\varepsilon}
% \newtheorem{assumption}{Assumption}
% \newtheorem{proposition}{Proposition}

%\documentclass{beamer}
\usepackage{beamerthemesplit} % new 
\usepackage{graphicx}
\usepackage{cases}
\usepackage{caption}
\usepackage{subcaption}
\usepackage{lmodern}
\usepackage{color}
\usetheme{Warsaw}

\title[FlowMkts]{Flow Markets}
\author[DF]{Daniel Friedman\thanks{Econ Departments, UCSC and U of Essex}
\and Yilin Li\thanks{Econ Department, UCSC}
\and Kristian Lopez Vargas\thanks{Econ Department, UCSC}
}
\useoutertheme{infolines}

\begin{document}

\date[UQ]{University of Queensland e-Seminar\\
August 28, 2024
} 

\frame{\titlepage} 


%\section{Introduction} 
% \title[FlowTrading]{Flow Trading Format: A Laboratory Test}
% \author[DF-YL-KL]{Dan Friedman, Yilin Li, Kristian L\'opez Vargas (UCSC)}
% \institute[UCSC]{}  

% \date[4/27/2024]{\vspace{1cm} \textbf{Material for UNSW Lecture 4} \\
% June 12, 2024}

% \begin{document}
% \maketitle

% \tableofcontents

%%%%%%%%%%%%%%%%%%%%%%%%%%%%%%%%%%%%%%
\section{Motivation}
\begin{frame}{
The continuous double auction (CDA) market format}
\begin{itemize}
\item dominated in 20th century financial markets 
\item but design is now problematic.  \pause
\end{itemize}
Problems arise from the conjunction of:
\begin{itemize}
\item strict price/time priority  
\item coarse price grid, and 
\item very low latency telecommunication and computing technology.

 \end{itemize} 
 \end{frame}

\begin{frame}{Creates a ``tilted playing field'': }
incentives to be the first to snap up stale bids/asks may lead to
 \pause
\begin{enumerate}
%\item Unequal opportunities to mitigate price impact
\item unproductive rent-seeking: 10s \$billions on sub-ms speed advantage \pause
\item reduced true liquidity due to toxic takers (``sniping'') \pause
\item increased danger of ``flash crashes,'' as algorithms interact. \pause
  \end{enumerate}
Also tilted by access to complicated ''order shredding'' algorithms.
\end{frame}

\begin{frame}{Reform proposals}
for financial market formats include:
    \begin{itemize}
\item IEX: see e.g., M. Lewis (2014), Aldrich and Friedman (2023); \pause
\item FBA: see e.g., Budish et al (2015), Aldrich and Lopez Vargas (2019); and \pause
\item Flow markets:  Kyle \& Lee (2017), Budish et al. (2023).
\end{itemize}
% The last is my personal favorite, and the subject of today's presentation.
\end{frame}

\begin{frame}{Flow market advocates claim}
that it will bring major benefits:
\begin{enumerate}
\item[0.] Implementable and understandable by human traders.
\item mitigates price impact 
(and volatility) by internalizing order shredding \pause
\item virtually eliminates sniping, and \pause
\item facilitates portfolio trades for interrelated assets.
\end{enumerate}
\end{frame}



\begin{frame}{Our project}

\textbf{compares CDA and Flow} market formats in a simple lab environment  with finely divisible asset units. 
\pause

This environment is 
suitable for testing claims 0 and 1.
\pause

It is \textbf{not} designed to test claims 2 and 3, though it does help set the stage. \pause 

\bigskip

A horse race (CDA vs Flow formats), not a test of theory.\\ \pause 
We compare across formats:
\begin{itemize}
\item traders' behavior \pause
\item market outcomes including 
\begin{itemize}
    \item price dynamics \pause
\item trade volume, and \pause
\item informational and allocative efficiency. 
\end{itemize} 
\end{itemize}


\end{frame}

\section{Market Formats}

\begin{frame}{Continuous Double Auction (CDA) Format}
processes standard limit orders.
\begin{itemize}
    \item These orders are defined by \textbf{three parameters}:$\rightarrow$ ``\underline{Buy/sell} up to \underline{$Q_{max}$} (shares) at a price of \underline{$P$} (dollars per share) or better."
\item Asynchronous order submission within a known time interval. 
\begin{itemize}
\item If a new order does not \textbf{cross} a queued contra-side order, then it is queued in the book. 
\item If it does cross, then it transacts immediately at the best price(s) on the other side of the market.
\end{itemize}
\item Strict \textbf{(price, time) priority}. If there is a tie on price, then older orders are filled first.
    \end{itemize} 
\end{frame}

\begin{frame}{CDA - Example}
    \begin{figure}[H]
        \centering
        \begin{tikzpicture}[scale=.6]
        \draw[thick,-] (0,0) -- (0,9) node[above]{$Price (\$/share)$};
        \draw[thick,-] (0,0)--(10,0) node[right]{$Quantity (shares)$};
        \node[left] at (0,1) {$41.21$};
        \node[left] at (0,2) {$41.22$};
        \node[left] at (0,3) {$41.23$};
        \node[left] at (0,4) {$41.24$};
        \node[left] at (0,5) {$41.25$};
        \node[left] at (0,6) {$41.26$};
        \node[left] at (0,7) {$41.27$};
        \node[left] at (0,8) {$41.28$};
        \node[left] at (0,9) {$41.29$};
        \node[below] at (2,0) {$500$};
        \node[below] at (4,0) {$1000$};
        \node[below] at (6,0) {$1500$};
        \node[below] at (8,0) {$2000$};
        \node[below] at (0,0) {$0$};
        \draw[blue,solid] (0.03,9) -- (0,3) -- (2,3) -- (2,3) -- (5,3) -- (5,2) -- (9,2) node[right]{Buy Orders} ;
        \draw[red,solid] (0.02,0) -- (0,4.05) -- (3.2,4.05) -- (3.2,5) -- (5,5) -- (5,6) -- (7,6) -- (7,8) -- (9,8) node[right]{Sell Orders};
        \node[above right] at (0,4.05) {Ask Price};
        \draw[fill=red] (0,4.05) circle [radius=.1];
        \node[below right] at (0,3) {Bid Price};
        \draw[fill=blue] (0,3) circle [radius=.1];
        \end{tikzpicture}
        \caption{Standard Limit Order Book}
    \end{figure}
\end{frame}


\begin{frame}{CDA - Example (at one instant)}
    \begin{figure}[H]
        \centering
        \begin{tikzpicture}[scale=.6]
        \draw[thick,-] (0,0) -- (0,9) node[above]{$Price (\$/share)$};
        \draw[thick,-] (0,0)--(10,0) node[right]{$Quantity (shares)$};
        \node[left] at (0,1) {$41.21$};
        \node[left] at (0,2) {$41.22$};
        \node[left] at (0,3) {$41.23$};
        \node[left] at (0,4) {$41.24$};
        \node[left] at (0,5) {$41.25$};
        \node[left] at (0,6) {$41.26$};
        \node[left] at (0,7) {$41.27$};
        \node[left] at (0,8) {$41.28$};
        \node[left] at (0,9) {$41.29$};
        \node[below] at (2,0) {$500$};
        \node[below] at (4,0) {$1000$};
        \node[below] at (6,0) {$1500$};
        \node[below] at (8,0) {$2000$};
        \node[below] at (0,0) {$0$};
        \draw[blue,solid] (0.03,9) -- (0,5) -- (2,5) -- (2,3) -- (5,3) -- (5,2) -- (9,2) node[right]{Buy Orders} ;
        \draw[red,solid] (0.02,0) -- (0,4.05) -- (3.2,4.05) -- (3.2,5) -- (5,5) -- (5,6) -- (7,6) -- (7,8) -- (9,8) node[right]{Sell Orders};
        \node[above right] at (0,4.05) {Ask Price};
        \draw[fill=red] (0,4.05) circle [radius=.1];
        \node[below right] at (0,3) {Bid Price};
        \draw[fill=blue] (0,3) circle [radius=.1];
        \end{tikzpicture}
        \caption{Standard Limit Order Book}
    \end{figure}
\end{frame}

\begin{frame}{FLOW Market Format}
processes continuous scaled limit orders (CSLOs).
    \begin{itemize}
\item A CSLO has \textbf{five parameters}: $\rightarrow$ ``\underline{Buy/Sell} up to a cumulative total of  \underline{$Q_{max}$} (shares) at maximum rate \underline{$U_{max}$} (shares per second) at prices between \underline{$P_{L}$} and \underline{$P_{H}$} (dollars per share)."
\item Consider a buy order. The (buying rate or) flow demand $U(p(t))$ is the decreasing pw linear function:
\begin{align}
                U{(p(t))}=
                \begin{cases}
                    U_{max} & \text{ if } p{(t)} < P_{L} \\
                    (\frac{P_{H}-p(t)}{P_{H}-P_{L}})U_{max} & \text{ if } P_L \leq p{(t)} \leq P_H \\
                    0  & \text{ if } p{(t)} > P_{H} \\
                \end{cases}
\end{align}
        where $p(t)$ is the clearing price at time $t$.
\item The flow supply (rate) is an analogous pw linear monotone \textbf{increasing} function of $p(t)$
    \end{itemize}
\end{frame}


\begin{frame}{FLOW format example - Individual}
    \begin{figure}
        \centering
        \includegraphics[width=\textwidth]{figures/flow.png}
        \caption{Continuous Scaled Limit Orders}
        \label{fig:flow}
    \end{figure}
\end{frame}


\begin{frame}{FLOW format example - Market}
    \begin{figure}[H]
        \centering
        \begin{tikzpicture}[scale=0.6]
            \draw[thick,-] (0,0) -- (0,9) node[above]{$Price (\$/share)$};
            \draw[thick,-] (0,0)--(9,0) node[right]{\begin{tabular}{c}
                 $Quantity \; Traded$ \\
                 $(1000 \; shares/hour)$
            \end{tabular}};
            \node[left] at (0,1.5) {$41.22$};
            \node[left] at (0,3.0) {$41.23$};
            \node[left] at (0,4.5) {$41.24$};
            \node[left] at (0,6.0) {$41.25$};
            \node[left] at (0,7.5) {$41.26$};
            \node[left] at (0,9.0) {$41.27$};
            \node[below] at (1,0) {$5$};
            \node[below] at (2,0) {$10$};
            \node[below] at (3,0) {$15$};
            \node[below] at (4,0) {$20$};
            \node[below] at (5,0) {$25$};
            \node[below] at (6,0) {$30$};
            \node[below] at (7,0) {$35$};
            \node[below] at (8,0) {$40$};
            \node[left] at (0,0) {$0$};
            \draw[blue,solid] (1,7.5) -- (2.5,6.0) -- (6.5,4.5) -- (7.5,3) -- (8,1.5) node[right]{Buy Orders} ;
            \draw[red,solid] (0.8,1.5) -- (1.8,3) -- (2,4.5) -- (5.5,6) -- (8,7.5) node[right]{Sell Orders};
            \draw[red,dashed] (2.5,6.0) -- (5.5,6.0);
            \draw[blue,dashed] (2,4.5) -- (6.5,4.5);
            \node[black,above] at (4,6.0) {Excess Supply};
            \node[black,below] at (4.25,4.5) {Excess Demand};
            \draw[dashed] (4.1,0) -- (4.1,5.4);
            \draw[dashed] (0,5.4) -- (4.1,5.4);
            \node[above right] at (4.1,0) {$U(t)=20.8$};
            \draw[fill] (4.1,0) circle [radius=.1];
            \node[left] at (0,5.4) {$P(t)=\$41.246$};
            \draw[fill] (0,5.4) circle [radius=.1];
        \end{tikzpicture}
        \caption{Continuous Scaled Limit Order Book}
    \end{figure}
\end{frame}

\begin{frame}{Flow format pricing}
Both the downward sloping demand schedule and upward sloping supply schedule are \textbf{piece-wise linear} functions of price $\rightarrow$ \textbf{unique} price and rate of shares per second. \pause
\begin{itemize}
\item All continuous scaled limit orders ares treated \textbf{symmetrically} and executed \textbf{simultaneously} at the market clearing price. \pause
\item Price changes only when a new order arrives, or an existing order achieves $Q_{max}$ or is cancelled.
\item Traders have \textbf{equal} access to \textbf{opportunities for order shredding}; $P_H - P_L$ and $U_{max}$ are  \textbf{``within-exchange''}
instructions to shred. \pause
\item Consequently, sophisticated algorithms to shred orders would seem unprofitable in the Flow format.
    \end{itemize} %\pause
%Aside: the ``portfolio'' version in Budish et al (2023) is a practical implementation of my monetary exchange process with unified prices. 
\end{frame}

\section{Lab Implementation}

\begin{frame}{Experiment Design: Main Elements}
    \begin{itemize}
        \item Markets of \textbf{8} traders, \textbf{20 paid periods} each (following 2 practice periods). \pause
        \item In each trading period each trader receives a \textbf{buy or sell contract} committing the experimenter to buy or sell up to $x$ shares at some point during the trading period 
        \item e.g., Computer will buy up to 300 shares from you at price 13 ECU at second 115) 
        \item Contracts are private information \pause
        \item simplifications: contracts all sent in first few seconds and are fulfilled in last few seconds; traders' inventories are restricted to $[-Q_{contract}, 0]$ for sellers or $[0, Q_{contract}]$ for buyers.
        \item Subjects can submit or update buy and/or sell orders at any time. \pause
        \item \textbf{Treatments}: \{CDA, FLOW\} (between-group)
        \item Three sets of contracts varied across 5 blocks of 4 periods each. 
    \end{itemize}
\end{frame}

\begin{frame}{Experiment Design: Contracts}
    
\begin{figure}
    \centering
    % \caption{Contract Configurations}
    \begin{subfigure}[c]{0.33\linewidth}
        \centering
            \begin{tikzpicture}[scale=.2]
                % create grid
                \draw[help lines, gray!30] (0,0) grid (18,20);
                
                % Draw axes
                \draw[->] (-1,0) -- (18,0) node[font=\tiny,right] {$Q$};
                \draw[->] (0,-1) -- (0,20) node[font=\tiny,above] {$P$};
              
                % Draw step functions
                \draw[blue, very thick] (0,20) -- (0,17) -- (6,17) -- (6,16) -- (9,16) -- (9,14) - - (13,14) -- (13,0) ;
                \draw[red, thick] (0,0) -- (0,3) -- (4,3) -- (4,4) -- (8,4) -- (8,11) --  (11,11) -- (11,14) -- (15,14) -- (15,20);
                
                % % Add labels
                \node[font=\tiny,below left] at (0,0) {$0$};
                \foreach \x in {5,10,15} {
                  \node[font=\tiny,below] at (\x,0) {$\x00$};
                }
                \foreach \y in {5,10,15} {
                  \node[font=\tiny,left] at (0,\y) {$\y$};
                }
                
                % % Add thicker ticks (Q)
                \draw[black, thick] (5,-0.3) -- (5,0.3);
                \draw[black, thick] (10,-0.3) -- (10,0.3);
                \draw[black, thick] (15,-0.3) -- (15,0.3);

                % % Add thicker ticks (P)
                \draw[black, thick] (-0.3,5) -- (0.3,5);
                \draw[black, thick] (-0.3,10) -- (0.3,10);
                \draw[black, thick] (-0.3,15) -- (0.3,15);
            \end{tikzpicture}
        \subcaption{\tiny Block 1 (T1-T4) \\ and Block 5 (T17-T20)}
    \end{subfigure}%
    \begin{subfigure}[c]{0.33\linewidth}
        \centering
        \begin{tikzpicture}[scale=.2]
            % create grid
            \draw[help lines, gray!30] (0,0) grid (18,20);
            
            % Draw axes
            \draw[->] (-1,0) -- (18,0) node[font=\tiny,right] {$Q$};
            \draw[->] (0,-1) -- (0,20) node[font=\tiny,above] {$P$};
          
            % Draw step functions
            \draw[blue, very thick] (0,20) -- (0,17) -- (4,17) -- (4,15) -- (8,15) -- (8,8) -- (12,8) -- (12,6) -- (16,6) -- (16,0);
            \draw[red, thick] (0,0) -- (0,3) -- (3,3) -- (3,4) -- (6,4) -- (6,5) -- (10,5) -- (10,6) -- (14,6) -- (14,20);
            
            % % Add labels
            \node[font=\tiny,below left] at (0,0) {$0$};
            \foreach \x in {5,10,15} {
              \node[font=\tiny,below] at (\x,0) {$\x00$};
            }
            \foreach \y in {5,10,15} {
              \node[font=\tiny,left] at (0,\y) {$\y$};
            }

            % % Add thicker ticks (Q)
            \draw[black, thick] (5,-0.3) -- (5,0.3);
            \draw[black, thick] (10,-0.3) -- (10,0.3);
            \draw[black, thick] (15,-0.3) -- (15,0.3);

            % % Add thicker ticks (P)
            \draw[black, thick] (-0.3,5) -- (0.3,5);
            \draw[black, thick] (-0.3,10) -- (0.3,10);
            \draw[black, thick] (-0.3,15) -- (0.3,15);
        \end{tikzpicture}
        \subcaption{\tiny Block 2 (T5-T8) \\ and Block 4 (T13-T16)}
    \end{subfigure}%
    \begin{subfigure}[c]{0.33\linewidth}
        \centering
        \begin{tikzpicture}[scale=.2]
            % create grid
            \draw[help lines, gray!30] (0,0) grid (18,20);
            
            % Draw axes
            \draw[->] (-1,0) -- (18,0) node[font=\tiny,right] {$Q$};
            \draw[->] (0,-1) -- (0,20) node[font=\tiny,above] {$P$};
          
            % Draw step functions
            \draw[blue, very thick] (0,20) -- (0,16) -- (3,16) -- (3,15) -- (7,15) -- (7,14) -- (11,14) -- (11,10) -- (15,10) -- (15,0);
            \draw[red, thick] (0,0) -- (0,3) -- (4,3) -- (4,4) -- (8,4) -- (8,5) -- (12,5) -- (12,8) -- (15,8) -- (15,20);
            
            % % Add labels
            \node[font=\tiny,below left] at (0,0) {$0$};
            \foreach \x in {5,10,15} {
              \node[font=\tiny,below] at (\x,0) {$\x00$};
            }
            \foreach \y in {5,10,15} {
              \node[font=\tiny,left] at (0,\y) {$\y$};
            }

            % % Add thicker ticks (Q)
            \draw[black, thick] (5,-0.3) -- (5,0.3);
            \draw[black, thick] (10,-0.3) -- (10,0.3);
            \draw[black, thick] (15,-0.3) -- (15,0.3);

            % % Add thicker ticks (P)
            \draw[black, thick] (-0.3,5) -- (0.3,5);
            \draw[black, thick] (-0.3,10) -- (0.3,10);
            \draw[black, thick] (-0.3,15) -- (0.3,15);
        \end{tikzpicture}
        \subcaption{\tiny Block 3 (T9 - T12)}
    \end{subfigure}
    \caption{\textbf{Contract configurations across trading blocks.}  \footnotesize
Vertical axis is price (ECU per share) and horizontal axis is quantity (shares). Each panel plots the step‐function demand schedule (thick blue) and supply schedule (red).  
The intersection of the two schedules determines the competitive‐equilibrium price–quantity pair for that block.  
Panel (a) applies to Blocks 1 (T1–T4) and 5 (T17–T20) of experiment; panel (b) to Blocks 2 (T5–T8) and 4 (T13–T16); panel (c) to Block 3 (T9–T12). }
% \label{fig:contract-configs}
    \label{fig:3contract_configs_new}
\end{figure}

\end{frame}


\begin{frame}{Experiment Design: Interface - CDA}
    \begin{figure}[H]
        \centering
        \includegraphics[width=\textwidth]{figures/cda_ui.png}
        \caption{CDA UI}
        \label{fig:cda_ui}
    \end{figure}
\end{frame}

\begin{frame}{Experiment Design: Interface - CDA Animation}
  \animategraphics[loop,controls,width=\linewidth]{10}{figures/cda_animation/cda_animation-}{0}{252}
\end{frame}



\begin{frame}{Experiment Design: Interface - FLOW Animation}
  \animategraphics[loop,controls,width=\linewidth]{10}{figures/flow_animation/flow_animation-}{0}{419}
\end{frame}


\begin{frame}{Experiment Design: Interface - FLOW}
    \begin{figure}[H]
        \centering
        \includegraphics[width=\textwidth]{figures/flow_ui.png}
        \caption{FLOW UI}
        \label{fig:flow_ui}
    \end{figure}
\end{frame}




\begin{frame}{Hypotheses}
    \begin{itemize}
\item[H1.] \textit{Traders will submit more but smaller orders in CDA than in Flow format.} \pause
\vspace{0.3cm}
        \item[H2.] \textit{Price volatility will be lower in Flow than in CDA format.}\pause
\vspace{0.3cm}        
        \item[H3.] \textit{Markets will be more liquid [lower PPI] in Flow than in CDA format.} \pause
\vspace{0.3cm}
        \item[H4.] \textit{The CDA and Flow formats will exhibit similar levels of trade volume.} \pause
\vspace{0.3cm}        
        \item[H5.] \textit{The CDA and Flow formats will exhibit similar levels of \\ (a) informational and (b) allocative efficiency.} 
    \end{itemize}
\end{frame}


\section{Results}

\begin{frame}{Summary of Results}

\begin{enumerate}
    \item[R1:] \textbf{Flow} Market format exhibits \textbf{fewer and larger orders}.  \pause 

    
    \item[R2:] \textbf{Flow} format exhibits \textbf{lower price volatility} compared to the CDA.\\  \pause 
   
     \item[R3:] \textbf{PPI} also considerably lower $\implies$ more liquid. \\  \pause 
    \vspace{0.5cm}
    
    \item[R4:] Trade \textbf{Volume is lower in the Flow} Market than in the CDA.  \pause 
    
    \item[R5:] Formats have \textbf{comparable} informational and allocative \textbf{efficiency}. \pause 
  \vspace{0.5cm}
    \item[R6:] (ex post) Buyers' excess profits are far larger than Sellers' in CDA, not so much in Flow. 
    
\end{enumerate}



\end{frame}


\begin{frame}{Descriptive Stats: Trader Behavior}
\begin{table}[H]
    \centering
    % \caption{Summary Statistics: Traders' Behavior}
    \resizebox{.95\textwidth}{!}{
    \begin{tabular}{lcccccc}
    \toprule
    ~                               & \multicolumn{2}{c}{T1 - T20} & \multicolumn{2}{c}{T1 - T10}  & \multicolumn{2}{c}{T11 - T20} \\
                                    & CDA       & FLOW      & CDA       & FLOW      & CDA       & FLOW      \\
    \midrule ↩
    % \multicolumn{7}{l}{\textbf{(a) Trader Behavior}} \\
    % \hline \\[-1.8ex]
     Number of Orders                       & 32       & 24            & 31       & 25            & 32        & 23            \\
     \hspace{3mm}Seller             & 17       & 12            & 17       & 13            & 17        & 11            \\
     \hspace{3mm}Buyer              & 15       & 12            & 15       & 12            & 15        & 12            \\
     % Order Size                     & 42       & 72            & 44       & 63            & 40        & 80            \\
     \\
     Quant. per Order                 & 113      & 182           & 115      & 180           & 111       & 185           \\
     \hspace{3mm}Seller             & 115      & 185           & 118      & 184           & 111       & 185           \\
     \hspace{3mm}Buyer              & 114      & 188           & 116      & 184           & 113       & 191           \\
     \\
     Order Price                    & 8.93     & [6.78, 12.64] & 9.18     & [7.15, 13.01] & 8.68      & [6.41, 12.28] \\
     \hspace{3mm}Seller             & 10.7     & [9.21, 15.46] & 11.13    & [9.76, 16.04] & 10.27     & [8.66, 14.88] \\
     \hspace{3mm}Buyer              & 6.9      & [4.22, 9.71]  & 7.01     & [4.26, 9.71]  & 6.79      & [4.19, 9.71]  \\

    \\
     Order Max Rate                 &           & 19.62         &           & 17.78         &           & 21.46         \\
     \hspace{3mm}Seller             &           & 19.67         &           & 18.06         &           & 21.29         \\
     \hspace{3mm}Buyer              &           & 20.15         &           & 17.93         &           & 22.37         \\
    \bottomrule
    \end{tabular}}
    \label{tab:summary_stat_trader}
\end{table}
\end{frame}

\begin{frame}{Descriptive: Transaction Price Dynamics}
    \begin{figure}[H]
        \centering
        \begin{subfigure}[b]{0.85\textwidth}
            \centering
            \includegraphics[width=1.15\linewidth]{figures/Grp-CDA-priceTrim.pdf}
            %\caption{CDA}
            \label{price_flow}
        \end{subfigure}
        \begin{subfigure}[b]{0.85\textwidth}
            \centering
            \includegraphics[width=1.15\linewidth]{figures/Grp-Flow-priceTrim.pdf}
           % \caption{FLOW}
            \label{price_cda}
        \end{subfigure}
       % \caption{Clearing Prices}
    \end{figure}
    % \vspace{-15pt}
    % \scriptsize{The vertical dotted lines are period breaks. The first 10 seconds from each trading period are left out. The lighter colors represent individual groups.} 
\end{frame}

\begin{frame}{Descriptive: Volume and Surplus Dynamics}
    \begin{figure}
        \centering
        % First row: Volume Trends
        \begin{subfigure}[b]{0.48\linewidth}
            \centering
            \includegraphics[width=\linewidth]{figures/groups_cda_quantity.png}
            \caption{CDA}
            \label{fig:combined_volume_cda}
        \end{subfigure}%
        \hfill % Space between the figures
        \begin{subfigure}[b]{0.48\linewidth}
            \centering
            \includegraphics[width=\linewidth]{figures/groups_flow_quantity.png}
            \caption{FLOW}
            \label{fig:combined_volume_flow}
        \end{subfigure}
        
        % Second row: Realized Surplus
        \begin{subfigure}[b]{0.48\linewidth}
            \centering
            \includegraphics[width=\linewidth]{figures/groups_cda_surplus.png}
            \caption{CDA - Realized Surplus}
            \label{fig:combined_surplus_cda}
        \end{subfigure}%
        \hfill % Space between the figures
        \begin{subfigure}[b]{0.48\linewidth}
            \centering
            \includegraphics[width=\linewidth]{figures/groups_flow_surplus.png}
            \caption{FLOW - Realized Surplus}
            \label{fig:combined_surplus_flow}
        \end{subfigure}
    \end{figure}
\end{frame}

\begin{frame}{Descriptive Stats: Market Performance}
    
    \begin{table}
    \centering
    % \caption{Summary Statistics: Market Performance}
    % \renewcommand{\arraystretch}{1.1}
     \resizebox{\textheight}{!}{
    \begin{threeparttable}
    \begin{tabular}{lcccccc}
    
    ~                               & \multicolumn{2}{c}{T1 - T20} & \multicolumn{2}{c}{T1 - T10}  & \multicolumn{2}{c}{T11 - T20} \\
                                    & CDA       & FLOW      & CDA       & FLOW      & CDA       & FLOW      \\
    \midrule 
     \\[-1.8ex] 
     \multicolumn{7}{l}{\textit{\textbf{(a)} Price Statistics} (Avg. Equil Price = $9.8$ ECUs)} \\
    % \hline 
     \\[-1.8ex]    
     Clearing Prices ($P_t$, ECUs)              & 8.62      & 9.37      & 9.04      & 9.57      & 8.21      & 9.18      \\
     \hspace{3mm}(Std. Dev.)                    & (2.77)    & (2.54)    & (2.95)    & (2.75)    & (2.51)    & (2.29)    \\
     $ \lvert  P_{t}-P_{t - 1} \rvert $         & 0.57      & 0.39      & 0.58      & 0.39      & 0.55      & 0.4       \\
     Std. Dev. of $(\ln P_{t} - \ln P_{t - 1})$ & 0.14      & 0.07      & 0.15      & 0.07      & 0.13      & 0.07      \\     
     Price Pressure Index (PPI)      & 3.98      & 0.89          & 4.09      & 1.02          & 3.88      & 0.76          \\
     \\ [-1.8ex]
     \multicolumn{7}{l}{\textit{\textbf{(b)} Volume, Transaction Intensity}} \\
     \\ [-1.8ex]
     Trade Volume (shares)                     & 1138      & 959       & 1097      & 872       & 1179      & 1046      \\
     Trade Vol. / $Q_{\text{CE}}$ (\%)        & 93.6      & 78.8      & 90.3      & 71.1      & 96.9      & 86.4      \\
     Filled Contract / $Q_{\text{CE}}$ (\%)     & 85.2      & 77.1      & 81.0      & 69.8      & 89.3      & 84.5      \\
     Num. of Transactions                             & 14        &           & 13        &           & 15        &           \\
     Mean Transaction Size                      & 78        &           & 82        &           & 74        &           \\
     \hspace{3mm}(Std. Dev.)             & (61)      &           & (61)      &           & (60)      &           \\
     Clearing Rate ($q/sec$)                 &           & 9.37      &           & 9.54      &           & 9.2       \\
     \hspace{3mm}(Std. Dev.)             &           & (2.6)     &           & (2.8)     &           & (2.4)     \\
      Time with no Flow Trans. (\%)               &           & 11          &           & 12          &           & 10           \\
     \\ [-1.8ex]
     \multicolumn{7}{l}{\textit{\textbf{(c)} Efficiency and Buyer-Seller Inequality}} \\
     \\ [-1.8ex]
     Info. Effic.: $\lvert P_t - P_{\text{CE}} \rvert $        & 2.16      & 2.08      & 2.19      & 2.32      & 2.14      & 1.84      \\     
     Alloc. Effic.: $\Pi$ / $\Pi_{\text{CE}}$ ($\times 100$)          & 84.3      & 81.81     & 79.76     & 77.31     & 88.85     & 86.31     \\
     
     ${\Pi^{\text{buy}}}/{\Pi^{\text{buy}}_{\text{CE}}} - {\Pi^{\text{sell}}}/{\Pi^{\text{sell}}_{\text{CE}}}$ 
     ($\times 100$) & 53.7      & 11.18     & 27.38     & -5.16     & 80.03     & 27.51     \\
     \\[-1.8ex] 

    \end{tabular}
    \begin{tablenotes}[para,flushleft]
      \footnotesize
    \end{tablenotes}
    \end{threeparttable}
   }
    \label{tab:summary_stat}
\end{table}
\end{frame}

\begin{frame}{Price Pressure Index (PPI)}
captures the per-share round trip  cost of buying a standard package ($Q^*=20$ shares) in a time interval ($\Delta t = 5$ sec). \pause
For CDA it is
\begin{equation} \label{eq:ppicda}
    PPI_{CDA} = \frac{1}{Q^*}\int_0^{Q^*} A(Q^{CE}+u,t) - B(Q^{CE}+u,t) du,
\end{equation}
the per-share area between orderbook S and D beyond the CE quantity, averaged over grid times $t=10, 15, ... , 110$ in each period. \pause
For Flow it is
\begin{equation} \label{eq:ppiflow}
    PPI_{Flow} = \frac{1}{Q^*}\int_t^{t+\Delta t} a(r+q,s) - b(r+q,s) ds,
\end{equation}
where $a$ and $b$ are flow S and D (at time $s$), while $r$ is the clearing flow rate (at intersection of $a$ and $b$) and 
$q=\frac{Q^*}{\Delta t}=4$ suffices to acquire the standard package over the time grid interval. Again we report averages over grid times $t=10, 15, ... , 110$ in each period. \pause

Caveat: calculation takes observed order books as given. Actual increases in buying and/or selling may provoke traders to alter their order streams.
\end{frame}



\subsection{Regression Analysis}

\begin{frame}{Regression Specs}
For each outcome metric, we estimate: 
\begin{equation} \label{eq:reg}
    y_{g,t} = \alpha + \sum^{4}_{i = 1} \left(\beta_i \; \text{Block}_i\right) + \gamma \; \text{FLOW}_{g,t} + \theta \; \text{Period}_{g,t} + \epsilon_{g,t}
\end{equation}
where $y_{g,t}$ is the observed outcome  for the given group (i.e., session) and time, and the $\gamma$ coefficient for dummy variable $\text{FLOW}_{g,t}$ is the treatment effect. \pause

We control for: 
\begin{itemize}
\item S-D configuration, via  block fixed effect coefficients $\beta_i$,  
\item time trends, via coefficient $\theta$ for period number $\text{Period}_{g,t}$, and
\item session effects, via clustered s.e.'s.
\end{itemize} \pause
 
% Dependent variables $y$ here are: deviations from CE $ =\mid P^W_t - P_{CE}\mid$, 
% mean absolute interval-to-interval price change $=\mid P^W_{t} - P^W_{t - 1}\mid$, and 
% volatility $=Std\left(\ln P^W_{t} - \ln P^W_{t - 1}\right).$
\end{frame}



%%%%%%%%%%%%%%%%%%%%%%
\begin{frame}{H1: Regressions for Number and size of orders}
\begin{table}[H]
\centering    
%\caption{Regression: Traders' Behavior}
\resizebox{1.2\textheight}{!}{
\begin{threeparttable}
\begin{tabular}{@{\extracolsep{5pt}}lcccc}
\toprule
& \multicolumn{2}{c}{Number of Orders} & \multicolumn{2}{c}{Order Size} \\
% & \cline{1-4}
\cmidrule{2-5}
& T1-T20 & T11-T20 & T1-T20 & T11-T20 \\
& (1) & (2) & (3) & (4) \\
\midrule
Intercept & 36.170$^{***}$ & 41.770$^{***}$ & 29.481$^{*}$ & 15.523$^{}$ \\
 & (4.702) & (6.710) & (15.726) & (23.527) \\
Flow Format (=1) & -7.840$^{***}$ & -8.820$^{***}$ & 10.314$^{*}$ & 17.437$^{***}$ \\
 & (1.512) & (0.827) & (5.319) & (4.700) \\
Period Number & -0.300$^{}$ & -0.576$^{*}$ & 2.384$^{***}$ & 2.946$^{**}$ \\
 & (0.219) & (0.347) & (0.805) & (1.296) \\
\midrule
Observations & 200 & 100 & 200 & 100 \\
Adjusted $R^2$ & 0.406 & 0.567 & 0.240 & 0.283 \\
\midrule
CDA avg. & 32 & 32 & 42 & 40 \\
\bottomrule
\end{tabular}
% \begin{tablenotes}[para,flushleft]
% \footnotesize
% \textit{Note:} $^{*}$p$<$0.1; $^{**}$p$<$0.05; $^{***}$p$<$0.01. The standard errors within parentheses are cluster-robust at the group level. We omit the coefficients of the block dummies. 
% \end{tablenotes}
\end{threeparttable}
}
\label{tab:TradersBehavior}
\end{table}
\end{frame}


\begin{frame}{Result 1: Flow Market format exhibits fewer and larger orders than CDA}
\begin{itemize}
\item Flow market format has 8-9 fewer orders submitted per period than the 32 in CDA ($p < 0.01$).
\item But the orders are larger, especially in later periods (by $\sim 17$, compared to $\sim$40 in CDA, ($p < 0.01$).  
\item Taken together, these results suggest that the Flow market format shreds orders more effectively than the CDA. 
% \item Overall lower trading volume in Flow, by $\sim$ 130 shares in later periods ($p < 0.05$).
\item Order size trends significantly upward over time.
\end{itemize}   
\end{frame}



%%%%%%%%%%%%%%%%%%%%%%%%

\begin{frame}{H2-H3: Price Volatility and PPI Regressions}

\begin{table}[H]
\centering    
%\caption{Regression: Price Volatility}
\resizebox{1.3\textheight}{!}{
\begin{threeparttable}
\begin{tabular}{@{\extracolsep{5pt}}lcccccc}
%\toprule
& \multicolumn{2}{c}{$\lvert P_{\tau} - P_{\tau - 1}\rvert$} & \multicolumn{2}{c}{SD$\left(\ln P_{\tau} - \ln P_{\tau - 1}\right)$} & \multicolumn{2}{c}{PPI_{\tau}} \\
\cmidrule{2-7}
& T1-T20 & T11-T20 & T1-T20 & T11-T20 & T1-T20 & T11-T20 \\
& (1) & (2) & (3) & (4) & (5) & (6) \\
\midrule
Intercept & 2.470$^{***}$ & 2.869$^{***}$ & 0.345$^{***}$ & 0.391$^{***}$ & 9.776$^{***}$ & 11.478$^{***}$ \\
  & (0.455) & (0.621) & (0.078) & (0.112) & (0.860) & (1.253) \\
Flow Format (=1) & -0.733$^{***}$ & -0.568$^{***}$ & -0.070$^{***}$ & -0.061$^{***}$ & -3.095$^{***}$ & -3.120$^{***}$ \\
  & (0.059) & (0.071) & (0.010) & (0.013) & (0.096) & (0.130) \\
Period number & -0.062$^{***}$ & -0.094$^{***}$ & -0.010$^{**}$ & -0.013$^{**}$ & -0.232$^{***}$ & -0.344$^{***}$ \\
  & (0.023) & (0.030) & (0.004) & (0.006) & (0.044) & (0.065) \\
Interval number & -0.002$^{}$ & 0.006$^{}$ & & & -0.100$^{***}$ & -0.070$^{***}$ \\
  & (0.004) & (0.004) & & & (0.009) & (0.011) \\
\midrule
Observations & 4,600 & 2,300 & 200 & 100 & 4,800 & 2,400 \\
Adjusted $R^2$ & 0.139 & 0.118 & 0.307 & 0.282 & 0.255 & 0.258 \\
\midrule
CDA avg. & 0.57 & 0.55 & 0.14 & 0.13 & 3.98 & 3.88 \\
\bottomrule
\end{tabular}
\begin{tablenotes}[para,flushleft]
\footnotesize
\textit{Note:} $^{*}$p$<$0.1; $^{**}$p$<$0.05; $^{***}$p$<$0.01. The standard errors within parentheses are cluster-robust at the group level. We omit the coefficients of the block dummies. 
\end{tablenotes}
\end{threeparttable}
}
\label{tab:PriceVolatility}
\end{table}

\end{frame}



\begin{frame}{Result 2: Flow format has lower price volatility than CDA}
\begin{itemize}
% \item \textbf{Price efficiency} is similar (insignificantly lower in Flow).
\item Significantly \textbf{smaller price changes} (by $\sim 0.7$, $p<0.01$) in Flow, 
\item Also, significantly lower \textbf{volatility} (by $\sim$7 pp, about half!, $p<0.01$).
\item Significant \textbf{improvement over time} in these metrics in both formats.
(Weak evidence elsewhere for more rapid improvement with Flow). 
\end{itemize}   
\end{frame}

\begin{frame}{Result 3: Flow format has far greater liquidity than CDA}
\begin{itemize}
\item significant increase in liquidity (reduction in PPI) within and across rounds.
\item PPI is reduced by $\sim 3.1 $ in Flow from $\sim$4 in CDA.
\item using 2 or 10 sec conventions for $\Delta t$ still cuts PPI to less than 1/3 of its CDA level.
\end{itemize}   
\end{frame}



%%%%%%%%%%%%%%%%%

\begin{frame}{H4: Trade Volume Regression}

\begin{table}
\centering    
%\caption{Regression: Traded Volume}
\resizebox{1.2\textheight}{!}{
\begin{threeparttable}
\begin{tabular}{@{\extracolsep{5pt}}lcccc}
%\toprule
& \multicolumn{2}{c}{Traded Volume} & \multicolumn{2}{c}{Filled Cont. Q / $Q_{CE}$} \\
\cmidrule{2-5}
& T1-T20 & T11-T20 & T1-T20 & T11-T20 \\
& (1) & (2) & (3) & (4) \\
\midrule
Intercept & 746.583$^{***}$ & 728.025$^{**}$ & 0.634$^{***}$ & 0.641$^{***}$ \\
  & (209.852) & (367.522) & (0.154) & (0.206) \\
Flow Format (=1) & -178.764$^{***}$ & -132.314$^{***}$ & -0.080$^{***}$ & -0.048$^{**}$ \\
  & (28.706) & (38.370) & (0.021) & (0.021) \\
Period number & 18.866$^{*}$ & 18.613$^{}$ & 0.016$^{*}$ & 0.014$^{}$ \\
  & (11.219) & (19.823) & (0.008) & (0.011) \\
\midrule
Observations & 200 & 100 & 200 & 100 \\
Adjusted $R^2$ & 0.454 & 0.322 & 0.264 & 0.150 \\
\midrule
CDA avg. & 1138 & 1179 & 0.852 & 0.893 \\
\bottomrule
\end{tabular}
\begin{tablenotes}[para]
\footnotesize
\textit{Note:} $^{*}$p$<$0.1; $^{**}$p$<$0.05; $^{***}$p$<$0.01. The standard errors within parentheses are cluster-robust at the group level. We omit the coefficients of the block dummies. 
\end{tablenotes}
\end{threeparttable}
}
\label{tab:Volume}
\end{table}

\end{frame}

\begin{frame}{Result 4: Trade Volume is lower in Flow format than in CDA.}
\begin{itemize}
% \item Flow market format has 8-9 fewer orders submitted per period than the 32 in CDA ($p < 0.01$).
% \item But the orders are larger, especially in later periods (by $\sim 17$, compared to $\sim$40 in CDA, $p < 0.01$).  
% \item Taken together, these results suggest that the Flow market format shreds orders more effectively than the CDA. 
\item Overall lower trading volume in Flow, still $\sim$ 130 shares in later periods ($p < 0.01$).
\item Volume in Flow trends upward over time (significantly so in specs not shown with trend*format interaction terms.).
\item Impact on filling contracts is somewhat less, but still significant at $p=.05$.
\end{itemize}   
\end{frame}





%%%%%%%%%%%%%%%%%%%




\begin{frame}{H5: Efficiency Regressions}

\begin{table}
\centering    
%\caption{Regression: Efficiency and Buyer-Seller Disparity}
\resizebox{1.35\textheight}{!}{
\begin{threeparttable}
\begin{tabular}{@{\extracolsep{5pt}}lcccccc}
%\toprule
& \multicolumn{2}{c}{$\lvert P_t - P_{\text{CE}}\rvert$} & \multicolumn{2}{c}{Realized Surplus (\%)} & \multicolumn{2}{c}{Buyer-Seller Disparity} \\
\cmidrule{2-7}
& T1-T20 & T11-T20 & T1-T20 & T11-T20 & T1-T20 & T11-T20 \\
%& (1) & (2) & (3) & (4) & (5) & (6) \\
\midrule
Intercept & 7.813$^{***}$ & 7.970$^{***}$ & 0.521$^{***}$ & 0.627$^{***}$ & 10327$^{***}$ & 14453$^{***}$ \\
 & (0.840) & (0.997) & (0.109) & (0.138) & (1809) & (2834) \\
Flow Format (=1) & -0.116$^{}$ & -0.191$^{}$ & -0.025$^{}$ & -0.025$^{}$ & -1608$^{}$ & -2082$^{*}$ \\
 & (0.218) & (0.328) & (0.035) & (0.030) & (1113) & (1072) \\
Period number & -0.228$^{***}$ & -0.234$^{***}$ & 0.020$^{***}$ & 0.014$^{**}$ & -92$^{}$ & -302$^{*}$ \\
 & (0.052) & (0.065) & (0.005) & (0.007) & (107) & (174) \\
\midrule
Observations & 4,800 & 2,400 & 200 & 100 & 200 & 100 \\
Adjusted $R^2$ & 0.317 & 0.446 & 0.291 & 0.052 & 0.779 & 0.805 \\
\midrule
CDA avg. & 2.16 & 2.14 & 84.3 & 88.85 & 2537 & 3417 \\
\bottomrule
\end{tabular}
\begin{tablenotes}[para]
\footnotesize
\textit{Note:} $^{*}$p$<$0.1; $^{**}$p$<$0.05; $^{***}$p$<$0.01. The standard errors within parentheses are cluster-robust at the group level. Buyer-Seller Disparity is defined as $({\Pi^{\text{buy}}}-{\Pi^{\text{buy}}_{\text{CE}}})-({\Pi^{\text{sell}}}-{\Pi^{\text{sell}}_{\text{CE}}})$. We omit the coefficients of the block dummies. 
\end{tablenotes}
\end{threeparttable}
}
\label{tab:EfficiencyAndDisparity}
\end{table}
\end{frame}


\begin{frame}{Result 5: Flow and CDA formats have comparable informational and allocative efficiency}
\begin{itemize}
\item Informational and allocative efficiencies differ insignificantly between formats.  

\item The deviation from $P_{CE}$ is lower in Flow, i.e., informational efficiency is better, but not sigificantly.

\item Allocative efficiency is 2.5\% worse in Flow, but not sigificantly. \pause
\vspace{1cm}
\item \textbf{Additional Result.} The Flow profit advantage (relative to CE) of buyers over sellers is less than half of that for CDA by the second half of periods, but it is so noisy that the treatment effect is significant only at 10\% level. 
\end{itemize}   

\end{frame}


\begin{frame}{Disparity in Buyer-Seller Profits}
    \begin{figure}[!htbp] 
    \centering
    \caption{Profit Distribution}
    \begin{subfigure}[b]{.5\linewidth}
        \centering
        \begin{adjustbox}{trim=0 0 0 {.11\height}, clip=true}
            \includegraphics[width=\linewidth]{figures/gross_profits_cdf.png}
        \end{adjustbox}
        \subcaption{\tiny Gross CDF}
    \end{subfigure}%
    \begin{subfigure}[b]{.5\linewidth}
        \centering
        \begin{adjustbox}{trim=0 0 0 {.11\height}, clip=true}
            \includegraphics[width=\linewidth]{figures/excess_profits_cdf.png}
        \end{adjustbox}
        \subcaption{\tiny Excess CDF}
    \end{subfigure}
    \label{fig:profits}
\end{figure}
\end{frame}




%%%%%%%%%%%%%%%%%%%%%%%%%%%%%%%%%%%%%%%%%%%%%%%%%%%%%%%%%%%%%%%%%%%%%%%%%%%%%%%%%%%%%


\section{Conclusions}

\begin{frame}{Conclusions}
\textbf{First empirical test} of Flow market format, comparing it to CDA in an induced-value setting with finely divisible units. \pause
\begin{itemize}
\item Flow starts worse than CDA but may improve faster (learning).
\item Within exchange \textbf{shredding} seems effective: fewer but larger orders in Flow than in CDA.
\item Price \textbf{volatility} is much lower in Flow. 
\item \textbf{Liquidity} is also much better (PPI lower). \pause
\item Some relative \textbf{undertrading} in Flow.
\item \textbf{Informational  $|P-P_{CE}|$ and allocative efficiencies} are comparable. \pause
\item Buyer-seller \textbf{asymmetry} in Flow seems reduced.\pause
\end{itemize} 
\end{frame}


\begin{frame}{Parting Remarks}

Flow works!
\begin{itemize}
    \item It is comprehensible by human traders, although perhaps they take longer to master it.
    \item It does well even in a simple contract environment that
\begin{itemize}
    \item aligns better with CDA limit orders, and
    \item offers no scope for portfolio trades or protection against predatory HFT traders.
\end{itemize}
\end{itemize}


\vspace{.4in}

\begin{center}
{\large A draft is available for those interested. \\ Suggestions are very welcome!}

\url{dan@ucsc.edu}
    
\end{center}


\end{frame}







%%%%%%%%%%%%%%%%%%%%%%%%%%%%%%%%%%%%%%%%%%%%%%%%%%%%%%%%%%%%%%%%%%%%%%%%%%%%%%%%%%%%%

\section{Appendix}


\begin{frame}{Continuous Double Auction (CDA) Format}
processes standard limit orders.
\begin{itemize}
    \item These orders are defined by \textbf{three parameters}:$\rightarrow$ ``\underline{Buy/sell} up to \underline{$Q_{max}$} (shares) at a price of \underline{$P$} (dollars per share) or better."
\item Asynchronous order submission within a known time interval. 
\begin{itemize}
\item If a new order does not \textbf{cross} a queued contra-side order, then it is queued in the book. 
\item If it does cross, then it transacts immediately at the best price(s) on the other side of the market.
\end{itemize}
\item Strict \textbf{(price, time) priority}. If there is a tie on price, then older orders are filled first.
    \end{itemize} 
\end{frame}



\begin{frame}{CDA Timing Issues}
1. Early bird special  
    \begin{itemize}
        \item Since market supply and demand are discontinuous step functions of price, there is \textbf{no guarantee of a single point of intersection}: there can be excess of demand (supply), and often there is. 
        \item Hence, \textbf{an allocation rule} to determine $\alpha(t_0, t)$ is needed (e.g., time priority, or pro-rata)
\item In most existing financial markets, orders are \textbf{prioritized} based on \textbf{the time they arrive}, so high speed is rewarded
\item If an agent can \textbf{react to a relevant signal} (even if public) \textbf{faster} (even if only a tiny bit) than other traders, then she can \textbf{pick off stale orders} from other traders and reverse position after the information has spread in the market.
\item  Thus: tilted playing field, rent-seeking, toxic takers $\implies$ illiquidity, increased flash crash risk.
    \end{itemize}
    2. To avoid adverse price pressure, traders will want to \textbf{shred} large orders.
    \end{frame}


\subsection{Experiment Implementation}

\begin{frame}{Procedures}
    \begin{itemize}
        \item We paid the cumulative profits of all 20 payment periods
        \item Exchange rate is 1000 ECUs = 1 USD and show-up fee is 7 USD \pause
        \item Experiment software was designed and developed at UCSC's LEEPS Laboratory
        \item Experiment sessions were conducted in person at LEEPS Laboratory
        \item So far have 5  FLOW and 5 CDA sessions
    \end{itemize}
\end{frame}

\subsection{Outcomes and Hypotheses}


\begin{frame}{Outcomes}
    \begin{itemize}
        \item For \textbf{trading behavior}, we study: 
     the number of orders, the size of the orders, submitted prices, and
     the maximum execution rate in Flow market. \pause
        \item For \textbf{market performance}, we study:
        \begin{itemize}
        \item \textit{Absolute Price Change}, $|P_t - P_{t-1}|$. 
        \item \textit{Price volatility}, $(\ln P_t - \ln P_{t-1})$ 
         \item \textit{Trading Volume}, $Q$ (at the market-period level)
\item\textit{Informational [in]Efficiency},  \( |P_t - P_{\text{CE}}| \). 
    \item\textit{Allocative Efficiency}, measured as \( \Pi/\Pi_{\text{CE}}\) (at the market-period level). 
        \end{itemize}
     
    \end{itemize}
\end{frame}


\end{document}